\scanpage{13}
\mspage{6}
We are going to introduce
\[
(M6)\quad\deleted{\unsure{\exists f_\alpha}}\quad
\left\{\begin{aligned}
&\forall \deleted{\unsure{X}}\inserted{Z},X\in\mathcal M_n,\quad Z\subset X,\quad\exists(f_\alpha)_{\alpha\in A}\ \text{finite family}\\
&\text{of tame functions }X\longrightarrow R\\
&\text{and some }(e_\alpha)_{\alpha\in A},\quad e_\alpha\in R_0,\quad\text{such that}\\
&Z=\bigcap f_\alpha^{-1}(e_\alpha).\quad\text{Moreover, we may assume that the}\\
&f_\alpha\text{ separate the points of }X\smallsetminus Z.
\end{aligned}\right.
\]

\margin{\uline{Proposition} $(M6)$, $(M7)$ imply the existence of the amalgamated sum with a free arrow.}
\margin{\diagram{\begin{tikzcd}[column sep=small]
X_1&Z\arrow[l,hook']\arrow[r]&X_2
\end{tikzcd}}}

\deleted{This implies that one can form a \inserted{cartesian} diagram}
\deleted{\diagram{\begin{tikzcd}
\deleted{X}\arrow[r,hook]&\deleted{R^n\times R^A}\\
\deleted{Z}\arrow[u,hook]\arrow[r,hook,no head]&\deleted{R^n\times\{e\}}\arrow[u,hook]
\end{tikzcd}}}
\deleted{$e=(e_\alpha)_{\alpha\in A}$}

\deleted{Let us apply this: $Z_1\hookrightarrow X_1$ and $Z_2\mathrel{\lhook\joinrel\relbar}X_2$}

which is \uline{equivalent} to the fact that the quotient $X/Z$ (contraction of $Z$ to a point) exists in $\mathcal E$ in the strict sense (i.e. $X\mathrel{\mapstochar\relbar}|X|$ commutes with it). Moreover, we assume
\margin{important \unsure{supplement})}
\[
(M7)\quad\forall Z,X\in\mathcal M_n,\quad Z\subset X,\quad\text{and every}
\]
\[
f:Z\longrightarrow R\text{ tame},\qquad\exists g:X\longrightarrow R\text{ tame}
\]
which extends \unsure{$f$}

\uline{Proposition} Assuming $M5,M6,M7$, the amalgamated sums $X_1\amalg_Z X_2$ ($Z\hookrightarrow X_1$, $Z\hookrightarrow X_2$) exist in $\operatorname{Esp}_{\mathcal M_*}$, and commute with the forgetful functor\ldots
\margin{cf generalization which comes below.}
\mspage{12}
